\section{A Certified, Attributed Evaluation Protocol}\label{sec:protocol}

Disprovers can be compared only if every tool is judged under the same semantics, the same admissibility rules and the same budget, and only if their output is judged for meaning as well as counted. The artifact repository lists all amendments and post-hoc deviations with timestamps, and Section~\ref{sec:threats} summarizes them.

\textbf{Independent certifier.} The certifier is a module of about 450 lines, plus a 300-line profile checker, and imports no generator code. Given a witness, which consists of rows per table, it executes the \emph{original} DDL on a fresh in-memory SQLite database and inserts the rows with parameter binding. It checks foreign keys at the value level, because some BIRD foreign keys reference non-unique columns and SQLite's enforcement therefore cannot be used. It then runs the \emph{original} gold and predicted strings without rewriting and evaluates the contracts and tiers of Section~\ref{sec:problem}. Every witness of every tool passes through the certifier, including the counterexamples released by SpotIt and SpotIt+. Lexical formats for $\mathcal A_{1f}$ come from the value samples that the SpotIt+ authors mined from the official BIRD databases, which cover 123 text columns.

\textbf{Populations.} The primary population consists of the EX-correct (prediction, gold) pairs. The secondary population, used for recall, consists of the EX-incorrect pairs whose prediction executes. Predictions that are identical to the gold query up to whitespace cannot differ, so they count as unresolved for every tool. Results on distinct (question, gold, prediction) triples are mapped back to every system that produced the triple.

\textbf{Budgets and acceptance.} Every self-run tool receives the same wall-clock budget per distinct pair on one core. Every tool's acceptance test is the certifier at the evaluated tier, so a tool that generates a non-admissible distinguishing database keeps searching until the budget expires. For SpotIt and SpotIt+, we use the authors' released per-question results, with bounds up to five, a median of 5.7\,s and a 95th percentile of about 20 minutes per pair. We read these results in two ways: the \emph{author-claimed} count of non-equivalent pairs, and the \emph{re-certified} count of released counterexamples that our certifier accepts. The first-stage hypotheses H1--H3 use the claimed counts, which is the reading most favorable to the baselines. In Stage~3, we also re-ran SpotIt from its released code under our budget (Section~\ref{sec:main}); the hypotheses SP1--SP2 compare certified counts.

\textbf{Attribution.}\label{sec:attribution} For each tool, we draw a random sample of 100 distinct certified pairs at its primary tier. The primary tier is $\mathcal A_0$ for the released SpotIt/SpotIt+ counterexamples, which rarely reach $\mathcal A_{1f}$, and $\mathcal A_{1f}$ or $\mathcal A_R$ otherwise. We classify the difference shown by each witness into the classes of Section~\ref{sec:meaningful}; prediction errors include pairs where both queries are wrong. The samples of all tools of one stage are pooled and shuffled, so annotators do not see which tool produced a witness. However, the blinding is imperfect. Synthetic values such as \texttt{\textquotesingle{}x1\textquotesingle{}} and \texttt{-1} are recognizable, and annotators could check whether a kind of data occurs by running read-only SQL against $D_0$. Every item is annotated twice independently under the same written instructions, and a third pass adjudicates disagreements. Annotation was performed by an AI model, Claude by Anthropic, which is also the model that implemented the tools. The first author checked two samples by hand (Section~\ref{sec:threats}). The estimated number of meaningful witnesses of a tool is its certified count times its sampled share. A sample of 100 pairs covers \V{attrQMin}--\V{attrQMax} distinct questions, so we report intervals from a bootstrap over questions.

\textbf{Held-out split and three stages.} The 1,534 BIRD development questions are split once by a salted hash. Three schemas are held out entirely. For the other eight schemas, 30\,\% of the questions, stratified by schema and difficulty, form the development set. Questions inspected in earlier phases of the project were forced into development, and eleven such questions of held-out schemas were excluded. In the \emph{first stage}, we developed methods on development questions only. We then fixed the protocol and hypotheses H1--H4, froze the code with its commit and SHA-256 hashes recorded, and ran the test questions once. The data-profile tier and CaliberGraph-R form the \emph{second stage}. They were conceived after we had seen the attribution of first-stage \emph{test} witnesses, so the idea of the tier and its constraint classes was informed by test data. For example, several classes, such as symmetric bonds and lookup tables, bind mainly on held-out schemas. The implementation and thresholds of the second stage were developed on development questions only. We then fixed hypotheses H5--H7, froze the code a second time and ran the test questions once. Section~\ref{sec:rresults} therefore also reports second-stage results separately for held-out and other schemas. A \emph{third stage} was pre-registered after both test runs and before any of its data were read. It tests the validity of the attribution and replicates the second stage on Spider, which no design decision had seen, and on thirteen unseen BIRD systems.

\textbf{Pre-registered hypotheses.} \emph{First stage.} \textbf{H1}: CaliberGraph certifies more EX-correct pairs at $\mathcal A_{1f}$ than SpotIt claims. \textbf{H2}: CaliberGraph+SC (CaliberGraph with the SpotIt+ constraints) certifies more at $\mathcal A_2$ than SpotIt+ claims. \textbf{H3}: CaliberGraph's recall on known positives exceeds those of SpotIt (claims) and ParSEval. \textbf{H4}: removing any one of anchoring, isolation, join-graph completion or the SQLite value domain significantly lowers the count. \emph{Second stage.} \textbf{H5}: CaliberGraph-R certifies more pairs at $\mathcal A_R$ than random fuzzing-R and ParSEval-R. \textbf{H6}: its estimated number of meaningful witnesses exceeds theirs. \textbf{H7}: its share of implausible witnesses is lower than CaliberGraph's at $\mathcal A_{1f}$. \emph{Third stage.} \textbf{A}: blind human annotators agree with each other and with the AI ($\kappa_2\ge0.6$) and find more implausible witnesses in the first stage than in the second. A is repeated in a separately pre-registered round~\textbf{R2} with facts of $D_0$. \textbf{B}: a second model agrees with the AI and shows the same stage difference. \textbf{C}: our issue list predicts the corrections of Jin et al.~\cite{jin2026annotation}. \textbf{S1}--\textbf{S3} repeat H5--H7 on Spider, and \textbf{E1} repeats H5 on unseen systems. Two further pre-registrations followed independent pre-submission reviews of the manuscript. \textbf{R3} is a third human round with three annotators and 150 witnesses. Its pre-registered outcomes are agreement $\ge0.6$ among the annotators and with the AI, the stage gap, and a different prediction-error share from the AI's. \textbf{SP} is a local re-run of SpotIt under our budget, with the hypothesis that CaliberGraph certifies more at $\mathcal A_{1f}$ and $\mathcal A_0$.

\textbf{What pre-registration means here.} Our pre-registrations are internally timestamped. Each is a document committed to the authors' git repository before the data that test its hypotheses were read. The pre-registrations of Stages 1--3 are timestamped by commits of our development repository, and the later ones also by the public artifact repository. No independent registry was used. The whole study, from the first pre-registration to the scoring of the third human round, took place within one week, from 13 to 19 October 2026. \ref{app:deviations} lists the deviations.

\textbf{Statistics.} Paired outcomes, i.e., witness or no witness, are compared with exact McNemar tests. Confidence intervals of count differences come from a bootstrap that resamples \emph{questions}, since the ten systems answer the same questions. Because McNemar tests treat pairs as independent, the question-level intervals are the primary evidence, and we print $p$-values below $10^{-10}$ as such. H1--H3 form a Holm-corrected family, with H3 treated as an intersection of two tests. The two comparisons of H5 form another such family. For H4, we add Holm-adjusted $p$-values over the four components post hoc. H6 was pre-registered with a parametric bootstrap from the attribution samples; we also report a bootstrap over questions. H7 is a one-sided two-proportion test.

\subsection{Use of AI models in this study}\label{sec:aiuse}
Following the journal's policy, we describe how AI was used in the research. The study was carried out with Claude Opus~5.5 (Anthropic; model \texttt{claude-opus-5-5}), which the authors operated in an agentic coding environment. Under the authors' direction, the model drafted the study design, the hypotheses and the pre-registrations, implemented all tools (CaliberGraph, CaliberGraph-R, the certifier, the profile miner and the baseline harnesses), ran all experiments and analyses, produced every AI attribution and root-cause label, and drafted the manuscript. The labels come from two independent passes with the prompt released in the artifact, followed by adjudication. Claude Sonnet~5.5 (\texttt{claude-sonnet-5-5}) served as the second annotator of the cross-model check (Section~\ref{sec:attrresults}).

The authors posed the research question, approved every design decision, pre-registration and code freeze, reviewed the code, the results and the text, and recruited the human annotators; the round-3 annotators reported working without AI tools, and for rounds 1--2 the authors state this (Section~\ref{sec:threats}). The authors take full responsibility for the content. Every number in the paper is generated by scripts from the saved outputs, and all prompts, code, outputs and annotations are released. Because the same model implemented the tools and labeled their witnesses, the attribution is not independent of the tool builder; Section~\ref{sec:attrresults} reports how far human annotators reproduce it, and Section~\ref{sec:threats} discusses the consequences.
